\documentclass[../main.tex]{subfiles}
\begin{document}

\chapter{Teorema chinezească a resturilor}

Teorema chinezească a resturilor (engl. \textit{The Chinese remainder theorem} sau \textit{CRT}) nu este strict necesară pentru olimpiade, dar ne ajută mult să ne familiarizăm cu algebra modulară

\begin{theorem}[Teorema chinezească a resturilor]
  Fie $n$ numere întregi $m_1, m_2, \dots, m_n$, astfel încît oricare două dintre ele sînt coprime. În plus, fie $n$ numere întregi $a_1, a_2, \dots, a_n$ astfel încît $0 \leq a_i < m_i, \forall i$. Atunci există un unic $x$ cu $0 \leq x < \prod_{i} m_i$ astfel încît $x \equiv a_i \pmod{m_i}, \forall i$.
\end{theorem}

Dacă sînteți familiari cu teoria numerelor, puteți citi direct:

\begin{itemize}
  \item \href{https://en.wikipedia.org/wiki/Chinese_remainder_theorem}{pagina Wikipedia};

  \item \href{https://sites.millersville.edu/bikenaga/number-theory/chinese-remainder-theorem/chinese-remainder-theorem.html}{o pagină} care tratează și cazul în care modulii nu sînt coprimi;

  \item \href{https://kconrad.math.uconn.edu/blurbs/ugradnumthy/crt.pdf}{o pagină} cu rezultate mai avansate (dincolo de scopul acestui capitol);

  \item pagina \href{https://cp-algorithms.com/algebra/chinese-remainder-theorem.html}{CP Algorithms} pentru cod.
\end{itemize}

Eu voi încerca să demonstrez lucrurile mai băbește și să arăt de unde scoatem unele formule (\textit{hint}: nu din inspirația divină).

\section{Un exemplu tabelar}

Iată tabelul complet pentru $m_1 = 4, m_2 = 7$.

\begin{table}[H]
  \centering
  \begin{tabular}{c|ccccccc}
    & $a_2 = 0$ & $a_2 = 1$ & $a_2 = 2$ & $a_2 = 3$ & $a_2 = 4$ & $a_2 = 5$ & $a_2 = 6$ \\
    \hline
    $a_1 = 0$ & 0 & 8 & 16 & 24 & 4 & 12 & 20 \\
    $a_1 = 1$ & 21 & 1 & 9 & 17 & 25 & 5 & 13 \\
    $a_1 = 2$ & 14 & 22 & 2 & 10 & 18 & 26 & 6 \\
    $a_1 = 3$ & 7 & 15 & 23 & 3 & 11 & 19 & 27 \\
    \hline
  \end{tabular}

  \caption{Toate valorile posibile pentru $a_1$ și $a_2$ cînd $m_1 = 4$ și $m_2 = 7$.}
\end{table}

Pornind de la acest tabel putem demonstra existența și unicitatea unei valori între 0 și 27 pentru fiecare pereche $(a_1, a_2)$. Într-adevăr, toate valorile între 0 și 27 își au locul undeva în tabel. Vom demonstra întîi că $x$-ul pentru fiecare celulă este unic. Să presupunem că două valori $x_1$ și $x_2$ ar sta în aceeași celulă. Atunci am avea că

$$
\begin{cases}
4 \ | \ (x_1 - x_2) \\
7 \ | \ (x_1 - x_2)
\end{cases}
$$

De aici deducem că

$$28 \ | \ (x_1 - x_2)$$

Singura explicație este că $x_1 = x_2$. Din unicitatea lui $x$ rezultă și existența lui $x$. Avem de distribuit 28 de numere în 28 de celule de capacitate 1, deci rezultă că toate celulele sînt ocupate. Acest argument se generalizează la oricîți moduli coprimi.

Acum să vedem cum îl putem afla algoritmic pe $x$.

\section{Un exemplu cu doi moduli}

Să rezolvăm sistemul de ecuații

$$
\begin{cases}
x &\equiv 3 \pmod{7}\\
x &\equiv 4 \pmod{11}
\end{cases}
$$

Cu puțin chin putem găsi răspunsul 59; de exemplu, putem enumera valorile $11k + 4$, respectiv 4, 15, 26, 37, 48, 59. Ne oprim pentru că $59 \equiv 3 \pmod{7}$.

Iată și o soluție constructivă. Dorim să scriem $x = 3P + 4Q$ și să găsim niște valori $P$ și $Q$ astfel încît

\begin{itemize}
  \item Modulo 7, $Q$ să dispară, iar $P$ să fie 1 (astfel rămînem cu 3).
  \item Modulo 11, $P$ să dispară, iar $Q$ să fie 1 (astfel rămînem cu 4).
\end{itemize}

Dacă vreți, am redus problema originală la două probleme similare, dar speciale: căutăm $P \equiv 1 \pmod{7}$ și $P \equiv 0 \pmod{11}$, respectiv $Q \equiv 1 \pmod{11}$ și $Q \equiv 0 \pmod{7}$.

Calculăm inversul fiecărui modul modulo celălalt (putem face asta, pentru că modulii sînt coprimi).

\begin{itemize}
  \item $8 \cdot 7 = 56 \equiv 1 \pmod{11}$, deci inversul lui 7 modulo 11 este 8.
  \item $2 \cdot 11 = 22 \equiv 1 \pmod{7}$, deci inversul lui 11 modulo 7 este 2.
\end{itemize}

Acum scriem $P = 2 \cdot 11$, tocmai pentru că el dispare modulo 11 și este 1 modulo 7. Similar, scriem $Q = 8 \cdot 7$, pentru că el dispare modulo 7 și este 1 modulo 11. Rezultă

$$x = 3 \cdot 2 \cdot 11 + 4 \cdot 8 \cdot 7 = 290$$

Este destul de clar că putem folosi orice valoare congruentă cu $x$ modulo 77, în particular $290 \mod 77 = 59$. Pentru cazul general, vezi Wikipedia.

\section{Un exemplu cu 3 moduli}

Să rezolvăm sistemul de ecuații

$$
\begin{cases}
x &\equiv 2 \pmod{5}\\
x &\equiv 6 \pmod{7}\\
x &\equiv 4 \pmod{11}
\end{cases}
$$

Ca mai sus, căutăm $x = 2P + 6Q + 4R$ unde:

\begin{itemize}
  \item Modulo 5, $Q$ și $R$ să dispară, iar $P$ să fie 1;
  \item Modulo 7, $P$ și $R$ să dispară, iar $Q$ să fie 1;
  \item Modulo 11, $P$ și $Q$ să dispară, iar $R$ să fie 1.
\end{itemize}

Rezultă că $P$ este multiplu de 7 și de 11 (deci de 77). Inversul lui 77 modulo 5 este 3. Deci fixăm $P = 7 \cdot 11 \cdot 3 = 231$. De amorul artei, am putea alege și valoarea $2 \cdot 77 = 154$, iar inversul acesteia ar fi 4. Deci am putea fixa și $P = 2 \cdot 7 \cdot 11 \cdot 4 = 616$. Dar nu am avea nimic de cîștigat.

Similar, inversul lui $5 \cdot 11$ modulo 7 este 6, deci $Q = 5 \cdot 11 \cdot 6$. În sfîrșit, inversul lui $5 \cdot 7$ modulo 11 este 6, deci $R = 5 \cdot 7 \cdot 6$. Punînd totul cap la cap, obținem:

$$
x = 2 \cdot (7 \cdot 11 \cdot 3) + 6 \cdot (5 \cdot 11 \cdot 6) + 4 \cdot (5 \cdot 7 \cdot 6) = 3282
$$

Ca mai sus, putem opera modulo $5 \cdot 7 \cdot 11 = 385$, așadar putem alege $x = 202$.

\section{Cazul general cu \texorpdfstring{$n$}{n} moduli}

Putem generaliza formula de mai sus observînd că $x$ este o sumă de termeni, cîte unul pentru fiecare rest și modul. Termenul pentru $a_i$ și $m_i$ este un produs dintre $a_i$, produsul tuturor modulilor în afară de $m_i$ și inversul acestui produs față de $m_i$. În limbaj matematic,

$$M \defeq \prod_{i=1}^{n} m_i$$

$$M_i \defeq \frac{M}{M_i}$$

$$b_i = M_i^{-1} \mod m_i$$

$$x = \sum_{i=1}^{n} a_i \cdot M_i \cdot b_i \mod M$$

Codul necesar este elementar:

\begin{minted}{c}
long long crt() {
  long long M = 1;
  for (int i = 0; i < n; i++) {
    M *= m[i];
  }

  long long x = 0;
  for (int i = 0; i < n; i++) {
    long long others = M / m[i];
    long long b = inverse(others, m[i]);
    x = (x + r[i] * others % M * b) % M;
  }

  return x;
}
\end{minted}

\section{Cazul cu moduli necoprimi}

Metoda de mai sus nu funcționează dacă modulii sînt necoprimi, deoarece poate eșua cînd încearcă să inverseze modulii unii față de alții (exemplu: 6 nu are invers față de 8). Studiem altă metodă mai generală, care desigur poate fi aplicată și în cazul anterior. Rețineți-o pe oricare vă place!

\subsection{Condiția de existență}

Mai întîi, condiția de existență a soluției. Să rezolvăm sistemul de ecuații

$$
\begin{cases}
x &\equiv 10 \pmod{28}\\
x &\equiv 11 \pmod{35}
\end{cases}
$$

Putem să aplicăm CRT în sens invers ca să desfacem modulii în valori coprime. Din fiecare ecuație de mai sus deducem două:

$$
\begin{cases}
x &\equiv 2 \pmod{4}\\
x &\equiv 3 \pmod{7}\\
x &\equiv 1 \pmod{5}\\
x &\equiv 4 \pmod{7}
\end{cases}
$$

Observăm imediat că sistemul nu are soluții datorită ecuațiilor (2) și (4). Descompunerea ne arată și cum putem repara contradicția. Pentru fiecare doi moduli necoprimi $m_1$ și $m_2$, condiția necesară pentru ca sistemul să aibă ecuații este

$$
\text{cmmdc}(m_1, m_2) \ | \ (a_1 - a_2)
$$

De exemplu, sistemul de mai jos are soluții pentru că $7 \ | \ (25-11)$:

$$
\begin{cases}
x &\equiv 25 \pmod{28}\\
x &\equiv 11 \pmod{35}
\end{cases}
$$

Dacă există, soluția este unică modulo $\text{cmmmc}(m_1, m_2, \dots, m_n)$ (în cazul de mai sus, 140).

\subsection{Metoda de construcție}

Metoda funcționează pentru două ecuații. Dacă avem trei sau mai multe, le procesăm pe primele două și obținem o nouă congruență modulo $\text{cmmmc}(m_1, m_2)$. La aceasta adăugăm a treia ecuație ș.a.m.d. Așadar, fie sistemul:

$$
\begin{cases}
x &\equiv a \pmod{m}\\
x &\equiv b \pmod{n}
\end{cases}
$$

Să scriem naiv ce înseamnă prima congruență. Există un $k$ astfel încît

$$x = km + a$$

Inserăm în a doua congruență:

$$
\begin{aligned}
km + a & \equiv b \pmod{n} \implies \\
km & \equiv (b-a) \pmod{n} \implies \\
\text{Există un $l$ astfel încît }\ km &= (b-a) + ln
\end{aligned}
$$

Ne-ar plăcea să spunem pur și simplu $k=(b-a) \cdot [m^{-1} \pmod{n}]$, dar inversul poate fi nedefinit. În schimb, fie $g = \text{cmmdc}(m, n)$. Să observăm în ultima ecuație că toți termenii se divid cu $g$, după cum rezultă din condiția de existență. Atunci scriem:

$$
\begin{aligned}
k \frac{m}{g} & = \frac{b-a}{g} + l \frac{n}{g} \\
k \frac{m}{g} & \equiv \frac{b-a}{g} \pmod{\frac{n}{g}} \\
k & \equiv \frac{b-a}{g} \left(\frac{m}{g}\right)^{-1} \pmod{\frac{n}{g}}
\end{aligned}
$$

De data aceasta inversul este definit, căci $m/g$ și $n/g$ sînt coprime și pozitive.

\end{document}
